\chapter{The song}\label{ch:song}

\section{Song lines}

The song area (the \texttt{SONG EDIT} part, \key{F4}) lists the song lines. Each column is a channel, and the number in a cell is the
track the channel plays on that line. The keys:

\begin{itemize}
  \item \key{0}--\key{F}, \key{CONTROL+LEFT}, \key{CONTROL+RIGHT} and \key{BACKSPACE} change the track number of the cell.
  \item \key{CONTROL+T} puts a new empty track in the cell, \key{CONTROL+D} makes a copy of the track and puts it in its place
        (the program asks when the track was used only once), \key{CONTROL+J} prepares a whole line of empty tracks.
  \item \key{INSERT}, \key{DELETE} (\key{CONTROL+I}, \key{CONTROL+U}) insert and delete song lines. The
        \texttt{GO TO LINE} values that point past the change are adjusted by the program.
  \item \key{CONTROL+K} inserts a copy or a clone of song lines, in the dialog of figure~\ref{fig:dialog-copy-lines}.
  \item \key{CONTROL+G} makes the line a \texttt{GO TO LINE}: the song continues from another song line, which loops it.
  \item \key{HOME} and \key{END}, \key{SHIFT+PAGE UP} and \key{SHIFT+PAGE DOWN} move to the start or the end of the song and
        between the subsongs.
\end{itemize}

\dialog{dialog-copy-lines}{Insert a copy or a clone of song lines.}{6.5cm}

\section{Subsongs}

A song may contain more than one tune: each subsong starts from a song line. The keys \key{SHIFT+PAGE UP} and
\key{SHIFT+PAGE DOWN} move between them, and the SAP export (chapter~\ref{ch:importexport}) lists them in its \emph{Subsongs} field: the
hexadecimal numbers of the song lines that start the subsongs, the first of which is the default one.

\section{Bookmarks}

\key{F8} sets a bookmark at the cursor and \key{CONTROL+F8} clears it. \key{SHIFT+F7} plays from the bookmark, a way to
listen to the same passage again and again while you work on it.

\section{The info area}

The \texttt{INFO EDIT} part (\key{SHIFT+F4}) edits the data of the song itself.

\begin{center}
\begin{tabular}{@{}lp{10cm}@{}}
\toprule
Field & Meaning \\ \midrule
\texttt{NAME} & the name of the song and of the author, up to 64 characters \\
\texttt{MUSIC SPEED: AA/MM/S} & \texttt{AA} is the current speed of the song, \texttt{MM} the main speed, which the song uses when it
   is played from the beginning (both from \texttt{\$01} to \texttt{\$FF}); \texttt{S} is the engine speed, from 1 to 8: how many times the
   player routine is called in each frame. A speed above 4 is not supported by the exported formats \\
\texttt{MAXTRACKLENGTH} & the maximal length of the tracks \\
\texttt{MONO / STEREO} & the number of tracks: 4 or 8 \\
\texttt{PAL / NTSC} & the video system, 50 or 60 frames per second \\
\bottomrule
\end{tabular}
\end{center}

The values are changed with the hexadecimal keys, or \key{CONTROL+arrow}. \key{TAB} and \key{LEFT}/\key{RIGHT} move between the
fields, \key{SHIFT+TAB} goes to the name and \key{ENTER} returns to the part you came from.

\section{Changing the kind of song}

The commands of the \menu{Song} menu that change the whole song open dialogs of their own.

\begin{itemize}
  \item \menu{Switch song between 4 or 8 channels\ldots} changes a mono song to stereo and the other way round.
        The data of the channels that disappear is lost, and the dialog says so.
  \item \menu{Song columns' order change/copy/clear\ldots} (figure~\ref{fig:dialog-columns}) reorders the columns of the song, copies
        one into another, turns a mono column into a stereo one or clears a column, for the song lines in a range.
  \item \menu{Change maximal length of tracks\ldots} (figure~\ref{fig:dialog-max-length}) changes the maximal length of the tracks.
  \item \menu{All size optimizations\ldots} runs the optimizations that make the song smaller: searching wise loops in all the
        tracks, clearing the unused ones, \ldots (section~\ref{sec:optimize}).
\end{itemize}

\dialog{dialog-columns}{The order of the columns of the song.}{8cm}
\dialog{dialog-max-length}{The maximal length of the tracks.}{7cm}

\section{Optimizing the size}\label{sec:optimize}

The Atari has little memory, so a song should be small. The program has the tools to make it so:
the \emph{wise loops} (a track that repeats a phrase is stored once), the clearing of the tracks and the instruments that
are not used in the song, the removal of duplicated tracks (the song is adjusted to the one left), and the
\emph{cleanup} commands of the \menu{Track} and \menu{Instrument} menus. The export for programmers (\emph{RMT stripped}) shows the
features used by the song, which lets a player routine be made smaller still (chapter~\ref{ch:importexport}).
